Dikaios Amanah is a Software Engineering capstone project that aims to bridge the financing gap faced by Small and Medium Enterprises (SMEs) in Bangladesh through a secure, transparent, and ethical digital investment platform. The proposed system is inspired by the Islamic finance principle of Musharakah, where investors and business owners share profits and losses instead of relying on interest-based lending.

Many SMEs struggle to obtain financing due to collateral requirements, limited access to venture capital, and a lack of trustworthy investment mechanisms. At the same time, many individual investors seek ethical investment opportunities that align with Islamic financial principles while offering greater transparency and accountability. Existing financing solutions rarely address the needs of both groups simultaneously.

Dikaios Amanah provides a digital platform that connects verified SMEs with verified retail investors through a structured investment process. Businesses submit funding requests, supporting documents, and business information, while investors browse verified investment opportunities and participate according to their financial capacity. The platform focuses on transparency, risk awareness, and secure information management rather than acting as a financial institution.

To enhance trust, the system incorporates a Permissioned Trust and Audit Layer that records critical events such as business verification, contract approvals, investment agreements, and profit declarations in a tamper-evident manner. Operational data is managed using a relational database, while cryptographic hashing and immutable audit records help maintain data integrity without replacing conventional database systems.

The proposed system will be implemented as a Minimum Viable Product (MVP) using modern web technologies including React, Node.js, Express.js, PostgreSQL, Firebase Storage, JWT authentication, bcrypt password hashing, and SHA-256 document hashing. Advanced features such as live payment gateway integration, regulatory compliance, and production-scale blockchain deployment are intentionally excluded from the MVP and considered future enhancements.
This project demonstrates the practical application of software engineering principles, secure system design, database management, authentication mechanisms, and ethical financial technology while addressing a real-world problem with a technically feasible solution.